\section{Stand der Forschung}
\label{sec:stand-der-forschung}

In der ersten Arbeit \cite{steidlehenssler2026milestone1} wurde eine Reihe von CNN-Architekturen für die Klassifikation von Szenen aus dem PatternNet-Datensatz untersucht.
Dabei standen vor allem der systematische Vergleich verschiedener Modellvarianten sowie der Einfluss von Architekturelementen wie Batch-Normalisierung, Dropout, Pooling und der Tiefe des Netzes im Mittelpunkt.
Die Ergebnisse dieser Vorarbeit lieferten eine belastbare CNN-Basis und zeigten, dass insbesondere batch-normalisierte Architekturen die beste Klassifikationsleistung erzielen.

Die vorliegende Arbeit baut auf dieser Grundlage auf und erweitert den bisherigen Ansatz um Knowledge Distillation.
Anstelle eines reinen Architekturvergleichs wird nun ein Teacher-Student-Setup betrachtet, in dem ein leistungsfähigeres Teacher-Modell Wissen an ein kompakteres Student-Modell überträgt.
Zusätzlich zur klassischen Auswertung der Test-Accuracy wird damit auch untersucht, wie sich der Wissenstransfer auf das Lernverhalten und das Leistungsniveau eines kleineren Modells auswirkt.

Im Vergleich zur ersten Arbeit liegt der zusätzliche Beitrag somit in der Distillation-Komponente selbst: Das Student-Modell soll nicht nur aus den harten Labels lernen, sondern auch von den weicheren Vorhersagen des Teachers profitieren.
Damit verschiebt sich der Fokus von der Frage, welche CNN-Architektur allein am besten abschneidet, hin zu der Frage, wie sich eine bereits gute Modellbasis durch Knowledge Distillation weiter verbessern und effizienter machen lässt.